\section{Give the sample stream a file representation}
The mixer and FIR routines worked on arrays of signed sample codes. Their arguments
supplied the addresses and lengths, so neither routine had to discover what an input
file meant. A WAV file adds that description: it identifies an audio representation,
channel count, sample rate and region of sample bytes. The file reader must establish
those facts before choosing which bytes to pass to a processing kernel.

In this chapter, read a bounded PCM WAV file, filter its audio and write a canonical
WAV result. Reuse the eight-tap Q15 arithmetic from the FIR chapter. The new work is
finding a valid sample slice, retaining its channel structure and carrying its history
across blocks. A successful file operation, a valid container and a correct filter are
separate observations. Each can fail while the other two appear plausible.

Begin with the four-frame input made by
\texttt{python3 drafts/pcm-wav-fixtures/fixture.py input.wav}. It is eighty-two bytes,
not a forty-four-byte header followed by audio. The same helper can create a two-second
input with \texttt{--fixture listen}. It refuses an existing output. Keep the short file
for exact tracing and the longer one for playback after the byte-level checks pass.
